\documentclass[11pt]{article}
\usepackage[margin=25mm]{geometry}
\usepackage{fontspec}
\setmainfont{Latin Modern Roman}
\usepackage{amsmath,amssymb,mathtools}
\usepackage{unicode-math}
\setmathfont{Latin Modern Math}
\usepackage{xurl,hyperref}
\hypersetup{colorlinks=true,urlcolor=blue}
\setcounter{secnumdepth}{0}
\setlength{\parindent}{0pt}
\setlength{\parskip}{5pt}
\setlength{\emergencystretch}{3em}
\newcommand{\sref}[2]{\hyperlink{source:#1:#2}{[S#2]}}
\newcommand{\cataloguescope}[1]{\paragraph{Recorded target.}#1}
\newcommand{\R}{\mathbb R}
\newcommand{\C}{\mathbb C}
\newcommand{\N}{\mathbb N}
\newcommand{\Z}{\mathbb Z}
\newcommand{\Q}{\mathbb Q}
\newcommand{\F}{\mathbb F}
\newcommand{\eps}{\varepsilon}
\begin{document}
% UnsolvedMath ID 30002147; source catalogue code OWR-12009-003
\setcounter{section}{30002146}
\section{Critical Ultrarelativistic Black-Hole Collision Solutions}\label{30002147}
{\small UnsolvedMath ID 30002147 · Mathematical Physics}
\subsection{Definitions and mathematical statement}

\paragraph{Shared notation.}$\mathbb N=\{1,2,\ldots\}$; the standard number fields have their usual meanings. All additional conventions are stated locally.


\paragraph{Definitions and conventions.}
The impact parameter is the transverse separation parameter of incoming equal-mass black holes in a scattering description. A critical impact parameter separates qualitatively different outcomes, such as scattering and merger. Nonspinning means zero initial intrinsic angular momentum; ultrarelativistic refers to large incoming boost or energy relative to rest mass. A helical Killing field combines time translation and rotation; an impulsive gravitational wave is a distributional spacetime idealization.

\paragraph{Target retained from the canonical source catalogue.}
\paragraph{Catalogue wording retained for traceability.}
For equal-mass nonspinning ultrarelativistic black-hole collisions near critical impact parameter, does a unique critical solution exist, and can its high-energy behavior be described by an adiabatic, helical-Killing, or colliding-impulsive-wave approximation?

\paragraph{Source scope and qualification.}
This is a numerical and analytical research programme. The source target leaves the limiting family, normalization and uniqueness notion to its collision model; no existence theorem for an unspecified class of initial data is asserted.
\subsection{Short English statement}
Is there a unique threshold geometry for high-energy black-hole collisions, and which proposed approximations describe it?
\subsection{Research attempt}
\paragraph{Model, date and investigated scope.}
OpenAI GPT-6 Astra (configured model \texttt{gpt-6-astra}), ultra reasoning; generated 11 October 2026. Derives a conditional no-radiation consequence of exact helical symmetry and the stated Bondi mass-loss law.

\paragraph{Sources and literature check.}
The complete supplied definition, including source qualifications and the dated literature assessment, was reviewed. The calculation below is an elementary restricted attempt motivated by [S2]; it does not claim novelty or a fresh literature-status determination.

\paragraph{Context and investigated scope.} The complete supplied definition, including source qualifications and the dated literature assessment, was reviewed. The calculation below is an elementary restricted attempt motivated by [S2]; it does not claim novelty or a fresh literature-status determination.

\paragraph{Testing exact helical symmetry against radiation loss.} Assume an asymptotically flat vacuum model with smooth future null infinity obeys the Bondi mass-loss identity
\[
 \frac{dM_B}{du}=-\frac1{32\pi}\int_{S^2}|N(u,\omega)|^2d\omega,
\]
where $u$ is retarded time and $N$ is the radiation news tensor. This identity is an explicit geometric input for the following conditional test; it is not derived here. Exact helical symmetry with nonzero angular velocity makes the sphere-integrated mass periodic in $u$ with period $T$, since a time advance by $T$ is a full spatial rotation.

\paragraph{A rigorous incompatibility within those assumptions.} Integrating over a period gives $0=M_B(u+T)-M_B(u)=-(32\pi)^{-1}\int_u^{u+T}\int|N|^2$. The integrand is nonnegative, so $N=0$ almost everywhere, and everywhere if smooth. Exact helical symmetry therefore cannot describe a persistently radiating isolated solution satisfying these assumptions. This does not prohibit an approximate helical description over a short interval or a solution with incoming radiation or different asymptotics.

\paragraph{Collision-target gap.} The source asks for existence and uniqueness of a critical collision limit, plus the validity of several high-energy approximations. No such limit is constructed. The calculation only restricts one approximation: an adiabatic or near-zone helical model must quantify its radiation-loss error rather than impose exact global symmetry. Proving compactness/selection of a near-critical family and controlling that error, or the impulsive-wave alternative, remains the first gap. Equal masses and initially zero spins do not supply these estimates by themselves.

\paragraph{Verification and adversarial audit.}
Checked periodic integrated mass and nonnegative news flux; explicitly kept the mass-loss law as an input and allowed approximate/local helical models. Model self-audit only; no independent human proof checking. Prompts A and B were applied to this canonical record and its qualified source scope. Checked periodic integrated mass and nonnegative news flux; explicitly kept the mass-loss law as an input and allowed approximate/local helical models. Retained unresolved status and zero original-target progress; No critical collision limit, uniqueness proof or controlled adiabatic/helical/impulsive approximation is constructed.

Coordinator checked the periodic mass-loss integral forcing vanishing news under the stated asymptotic identity. This additional check was model-assisted, not independent human verification.

\paragraph{Tools, process and contributions.}
Python 3 was used for file preparation and structural checks; no numerical experiment or theorem prover was used. Contribution prompt A was followed by prompt B. GPT-6 Astra with ultra reasoning drafted the research and performed its stated self-audit; any additional coordinator contribution is identified in the verification and audit. A GPT-6-based Codex coordinator assembled the manuscript and checked repository integration; its exact variant and reasoning configuration were not exposed. No human mathematical contribution, independent human proof check or formal proof certification is supplied. Complete derivations and the consulted source references are included here; private reasoning transcripts are not published.

\paragraph{Final status and first remaining gap.}
\textbf{UNRESOLVED}. No critical collision limit, uniqueness proof or controlled adiabatic/helical/impulsive approximation is constructed.

\section{Completion Estimate}
\noindent\textbf{COMPLETION: 0\%}
The calculation is an auxiliary or restricted attempt and establishes no newly verified portion of the full source target. No critical collision limit, uniqueness proof or controlled adiabatic/helical/impulsive approximation is constructed.

\subsection{Sources}
\begin{itemize}
\item[\hypertarget{source:30002147:1}{S1}] ULAM.AI and the UnsolvedMath Contributors, UnsolvedMath, record 30002147 (OWR-12009-003). Catalogue curation under CC BY 4.0; definitions and English exposition added locally.
\url{https://huggingface.co/datasets/ulamai/UnsolvedMath}.
\item[\hypertarget{source:30002147:2}{S2}] Mathematical Aspects of General Relativity (2012), Oberwolfach Reports, DOI 10.4171/owr/2012/37, pp. 2269--2333. Primary report accessed 2026-10-11; relevant PDF pages 31, 32, 33.
\url{https://doi.org/10.4171/owr/2012/37}.
\item[\hypertarget{source:30002147:3}{S3}] Klein, helical binary spacetimes (catalogue-listed context; not independently consulted in this preparation).
\url{https://arxiv.org/abs/gr-qc/0410095}.
\item[\hypertarget{source:30002147:4}{S4}] Healy--Ciarfella--Lousto, high-speed collisions (catalogue-listed context; not independently consulted in this preparation).
\url{https://doi.org/10.1103/PhysRevD.111.084077}.
\end{itemize}
\label{end:30002147}
\end{document}
